% !TEX root = ../../ctfp-print.tex

\lettrine[lhang=0.17]{如}{果我还没有}说服你范畴论的核心是态射(morphism)，那我的工作就没有做到位。由于下一主题是伴随性(adjunctions)——通过hom集的同构定义的概念——回顾我们对hom集基本构成的理解是有必要的。此外你会发现伴随性提供了一种更通用的语言来描述许多之前研究过的构造，因此复习相关内容也会有所帮助。

\section{函子}

首先你应该真正把函子(functor)视为态射的映射——这正是Haskell中\code{Functor}类型类定义所强调的观点，该定义围绕着\code{fmap}展开。当然函子也映射对象——即态射的端点——否则我们就无法讨论复合性的保持。对象指明了哪些态射对是可组合的：若两个态射要进行组合，则一个态射的目标必须等于另一个态射的源。因此如果我们希望态射的复合对应到提升态射(lifted morphism)的复合，它们端点的映射方式基本上就确定了。

\section{交换图}

许多态射的性质都是通过交换图来表达的。如果某个特定的态射能够以多于一种的方式表示为其他态射的复合，则我们得到一个交换图。

特别地，交换图构成了几乎所有普遍构造（universal construction）的基础（显著例外是初始对象（initial object）和终止对象（terminal object））。我们在乘积（product）、余乘积（coproduct）、各种（上）极限（(co-)limit）、指数对象（exponential object）、自由幺半群（free monoid）等概念中都已见到过这种构造。

乘积是一个简单的普遍构造示例。我们选取两个对象 $a$ 和 $b$，并考察是否存在一个对象 $c$，以及一对态射 $p$ 和 $q$，使得该结构具有作为其乘积的泛性质（universal property）。

\begin{figure}[H]
  \centering
  \includegraphics[width=0.3\textwidth]{images/productranking.jpg}
\end{figure}

\noindent
乘积是极限（limit）的一个特例。极限是通过锥（cone）来定义的。一般性的锥由交换图构成。这些交换图的交换性可以用函子映射的适当自然性条件（naturality condition）替代。通过这种方式，交换性被降级为自然变换这一更高阶语言的汇编语言。

\section{自然变换}

通常情况下，当我们需要从态射到交换方块的映射时，自然变换是非常便捷的工具。一个自然性方块的两条对边分别表示某个态射 $f$ 在两个函子 $F$ 和 $G$ 下的映射结果。另外两边则是自然变换的分量（这些分量本身也是态射）。

\begin{figure}[H]
  \centering
  \includegraphics[width=0.35\textwidth]{images/3_naturality.jpg}
\end{figure}

\noindent
自然性意味着当你移动到"相邻"分量时（这里所说的相邻是指通过某个态射连接），你不会违背范畴或函子的结构。无论你是先使用自然变换的某个分量在对象之间建立桥梁，再通过函子跳转到其相邻对象；还是采用相反的操作顺序，最终结果都是相同的。这两种操作方向是正交的——可以形象地说，自然变换让你在水平方向移动，而函子则负责垂直方向或前后方向的移动。你可以将函子的\emph{像}可视化为目标范畴中的一个平面层。自然变换则将对应于 $F$ 的这种平面层映射到对应于 $G$ 的另一平面层。

\begin{figure}[H]
  \centering
  \includegraphics[width=0.35\textwidth]{images/sheets.png}
\end{figure}

\noindent
我们在Haskell语言中见过这种正交性的例子。在那里，函子的作用是在不改变容器形状的前提下修改容器内容，而自然变换则将未改变的内容重新打包到另一个容器中。这两个操作的顺序并不影响最终结果。

我们曾看到在极限定义中的锥体（cone）可以被自然变换替代。自然性保证了每个锥体的所有边都满足交换性。然而，极限本质上是通过锥体之间的映射来定义的。这些映射也必须满足交换性条件。（例如，在积的定义中涉及的三角形必须交换）

这些条件同样可以用自然性来替代。你可能还记得，\emph{泛}锥体（即极限）被定义为两个（逆变）Hom函子之间的自然变换：
$$F \Colon c \to \cat{C}(c, \Lim[D])$$
与另一个（同样是逆变的）函子之间的自然变换：
$$G \Colon c \to \mathit{Nat}(\Delta_c, D)$$
其中 $\Delta_c$ 是常值函子，$D$ 是定义范畴 $\cat{C}$ 中图示的函子。函子 $F$ 和 $G$ 都在 $\cat{C}$ 的态射上有明确的作用方式。特别地，这个特定的 $F$ 与 $G$ 之间的自然变换实际上是一个\emph{同构}。

\section{Natural Isomorphisms}

自然同构（natural isomorphism）——这是指其每个分量均可逆的自然变换——是范畴论表达「两件事物相同」的方式。此类变换的每个分量都必须是对象间的同构（isomorphism）——即具有逆态射（inverse morphism）的态射。若将函子像（functor images）可视化为层状结构，则自然同构就是这些层之间的逐点可逆映射。

\section{同态集合（Hom-Sets）}

但什么是态射（morphisms）呢？它们确实比对象具有更多的结构：与对象不同，态射有两个端点。但如果固定源对象和目标对象，两者之间的态射就构成了一个平凡的集合（至少对局部小范畴而言）。我们可以为这个集合中的元素起名为$f$或$g$以示区分——但究竟是什么使它们真正不同呢？

给定同态集合（hom-set）中态射的本质区别在于它们如何与其他态射（来自相邻同态集合）进行组合。若存在某个态射$h$，它与$f$的组合（无论是前组合还是后组合）不同于与$g$的组合，例如：
$$h \circ f \neq h \circ g$$
那么我们就可以直接"观测"到$f$与$g$之间的差异。即使这种差异无法被直接观测到，我们也可以通过函子（functor）放大同态集合的细节。某个函子$F$可能会将这两个态射映射为更丰富范畴中的不同态射：
$$F f \neq F g$$
此时相邻同态集合提供了更多可分辨的细节，例如存在某个不在$F$像集中的$h'$满足：
$$h' \circ F f \neq h' \circ F g$$

\section{同态集同构}

许多范畴论构造依赖于同态集之间的同构关系。但由于同态集本质上只是集合，单纯的集合间同构并不能提供太多信息。对于有限集合来说，同构仅表示它们包含相同数量的元素；若集合无限，则需具有相同的基数。但任何有意义的同态集同构必须考虑复合运算。而复合涉及多个同态集之间的交互。我们需要定义跨越整个同态集族的同构，并施加与复合运算兼容的相容性条件。此时，\newterm{自然}同构恰好能满足这一需求。

但什么是同态集的自然同构呢？自然性的本质是函子间的映射性质，而非集合间的性质。因此我们实际上讨论的是取值于同态集的函子间的自然同构。这些函子不仅仅是集合赋值函子，其在态射上的作用由相应的hom函子通过前复合或后复合（取决于函子的协变性）诱导而来。

Yoneda嵌入就是这种同构的一个典型例子。它将范畴$\cat{C}$中的同态集映射到函子范畴中的同态集，并且具有自然性。Yoneda嵌入中的一个函子是$\cat{C}$中的hom函子，另一个则将对象映射到同态集之间的自然变换集合。

极限的定义同样是一个同态集间的自然同构（第二个同构仍在函子范畴中）：
$$\cat{C}(c, \Lim[D]) \simeq \mathit{Nat}(\Delta_c, D)$$
事实上，我们之前构造的指数对象或自由幺半群也可以重新表述为同态集间的自然同构。

这绝非偶然——接下来我们将看到，这些不过是伴随关系(adjunction)的不同实例，而伴随关系正是通过同态集的自然同构来定义的。

\section{Hom集的不对称性}

还有一个观察结果有助于我们理解伴随关系(adjunctions)。
一般来说，Hom集并不具有对称性。一个Hom集$\cat{C}(a, b)$往往与反向的Hom集$\cat{C}(b, a)$有显著差异。
这种不对称性的极端示例是将偏序(partial order)视为范畴的情形。
在偏序范畴中，从$a$到$b$的态射存在当且仅当$a \leq b$成立。
若$a$与$b$不同，则不可能存在反向态射（即从$b$到$a$的态射）。
因此当Hom集$\cat{C}(a, b)$非空时（在此情形下必为单元素集合），除非$a = b$，否则$\cat{C}(b, a)$必定为空集。
该范畴中的箭头具有明确的单向流动特性。

基于非反对称关系的预序(preorder)，也具有"基本定向"特征，仅存在少量环状例外。
将任意范畴视为预序的推广形式是一种便捷的理解方式。

预序是一个薄范畴(thin category)--- 所有Hom集要么是单元素集合，要么为空集。
我们可以将一般范畴可视化为"厚化"的预序。

\section{挑战}

\begin{enumerate}
  \tightlist
  \item
        考虑自然性条件的一些退化情况，并绘制相应的示意图。例如，当函子$F$或$G$将$f \Colon a \to b$的定义域与陪域对象$a$和$b$映射到同一对象时（如$F a = F b$或$G a = G b$），会发生什么现象？
        （注意通过这种方式会得到一个锥体(cone)或余锥体(co-cone)。）接着考虑$F a = G a$或$F b = G b$的情形。
        最后，若从自环态射$f \Colon a \to a$出发又会产生何种情况？
\end{enumerate}